\documentclass[11pt,a4paper]{article}

\usepackage{amsmath,amssymb,amsthm}
\usepackage{booktabs}
\usepackage{hyperref}
\usepackage{xcolor}
\usepackage[margin=25mm]{geometry}
\usepackage{xeCJK}
\setCJKmainfont{Hiragino Mincho ProN}
\setCJKsansfont{Hiragino Sans}

\newtheorem{theorem}{定理}
\newtheorem{proposition}{命題}
\newtheorem{corollary}{系}
\newtheorem{observation}{観察}

\title{算術的ランドスケープ:\\
$L$関数による真空選択とダークエネルギーの情報コスト}

\author{Wright Brothers Research Program}
\date{2026年4月}

\begin{document}
\maketitle

\begin{abstract}
我々は，可能な真空状態の空間がCalabi--Yau多様体によるコンパクト化
（弦理論ランドスケープ，$\sim 10^{500}$個の真空）ではなく，
Dirichlet $L$関数によって分類されることを提案する．
各$L$関数 $L(s, \chi)$ は，$s = -1$ で正則化された零点エネルギーを
通じて真空を定義する:
$\Omega_\Lambda(\chi) = (8\pi/3)|L(-1, \chi)|$．
物理的制約 ($\Omega_\Lambda > 0$, $\Omega_\Lambda < 1$) が
この\emph{算術的ランドスケープ}を\textbf{本質的にただ一つの真空}に
帰着させることを証明する:
$L(s, \chi_0 \bmod 2) = \zeta_{\neg 2}(s)$，
$\Omega_\Lambda = 2\pi/9$ を与える．
具体的には，導手 $q \leq 100$ の全ての原始的Dirichlet指標のうち，
$q = 2^k$ のもののみが物理的真空を与え，全て同一の $\zeta_{\neg 2}$ に帰着する．
さらに，``標準的''真空 (DD境界, $\zeta$) から物理的真空
(DN境界, $\zeta_{\neg 2}$) への遷移が，
正確に $\Delta S = \frac{1}{2}\log 2$ nats --- 半ビット --- の
情報論的コストを持つことを示す．
これはAffleck--Ludwig境界エントロピー差に対応する．
宇宙論定数は特定の\textbf{情報論的}量に結び付けられる:
ダークエネルギーは宇宙の境界非対称性の熱力学的コストである．
算術的ランドスケープは有限 ($\leq 6$個の真空) かつ実質的に一意であり，
Euclid/DESIにより $0.5\%$ の精度で反証可能である．
\end{abstract}

\section{序論}

\subsection{ランドスケープ問題}

基礎物理学の中心的な課題は真空選択である:
なぜ宇宙はこの特定の真空状態を占めているのか?
弦理論は $\sim 10^{500}$ 個の準安定真空
（``ランドスケープ''）を予言するが~\cite{Susskind2003,KKLT2003}，
観測されたものを選択する既知の原理は存在しない．
宇宙論定数 $\Lambda$ は異なる真空で異なる値をとり，
その観測された小ささ ($\Omega_\Lambda \approx 0.685$) は
人間原理的選択に帰される~\cite{Weinberg1987}．

本論文は代替案を提案する: 真空の空間は幾何学的コンパクト化ではなく
\textbf{Dirichlet $L$関数}によって分類され，物理的制約が
この\emph{算術的ランドスケープ}を本質的にただ一つの真空に帰着させる．

\subsection{Euler積から真空状態へ}

以前の研究~\cite{wb-functional}において，我々は真空零点エネルギーが
Riemann $\zeta$ 関数のEuler積を通じて境界条件に依存することを示した．
Dirichlet--Neumann (DN) 境界条件は
$\zeta_{\neg 2}(s) = (1 - 2^{-s})\zeta(s)$ を選択し，
宇宙論定数 $\Omega_\Lambda = 2\pi/9$ は
時間的DN境界を持つWheeler--DeWitt方程式から
導かれる~\cite{wb-wdw}．

ここで一般化を行う: \emph{任意の} Dirichlet指標 $\chi$ が
$L(s, \chi)$ を通じて真空を定義し，
物理的制約 ($0 < \Omega_\Lambda < 1$) が\textbf{選択原理}として機能する．

\subsection{情報論的接続}

DD境界条件からDN境界条件への遷移は，真空のエンタングルメント
エントロピーを正確に $\Delta S = \frac{1}{2}\log 2$ だけ変化させる．
これは\textbf{Affleck--Ludwig境界エントロピー}~\cite{AffleckLudwig1991}
--- 境界共形場理論における普遍的な量 --- である．
従って，ダークエネルギーは情報論的解釈を獲得する:
宇宙の時間的境界非対称性 (ビッグバンにおけるDirichlet条件，
ド・ジッター未来におけるNeumann条件) の
\emph{熱力学的コスト}である．

\section{算術的ランドスケープ}

\subsection{真空状態としてのDirichlet $L$関数}

法 $q$ のDirichlet指標 $\chi$ は，周期 $q$ の完全乗法的関数
$\chi: \mathbb{Z} \to \mathbb{C}$ である．
対応する $L$関数は
\begin{equation}
L(s, \chi) = \sum_{n=1}^\infty \frac{\chi(n)}{n^s}
= \prod_{p\text{ 素数}} \frac{1}{1 - \chi(p)p^{-s}}.
\end{equation}

法 $q$ の主指標 $\chi_0$ に対して:
\begin{equation}
L(s, \chi_0) = \zeta(s)\prod_{p | q}(1 - p^{-s}).
\end{equation}

各 $L$関数は $L(-1, \chi)$ を通じて正則化されたエネルギーを持つ
``真空''を定義する:
\begin{equation}\label{eq:omega-chi}
\Omega_\Lambda(\chi) = \frac{8\pi}{3}|L(-1, \chi)|.
\end{equation}

値 $L(-1, \chi)$ は一般化Bernoulli数から計算される:
\begin{equation}
L(-1, \chi) = -\frac{B_{2,\chi}}{2}, \quad
B_{n,\chi} = q^{n-1}\sum_{a=1}^{q}\chi(a)B_n(a/q)
\end{equation}
ここで $B_n(x)$ は第 $n$ Bernoulli多項式である．

\subsection{真空の分類}

導手 $q \leq 6$ の全てのDirichlet指標に対して
$\Omega_\Lambda(\chi)$ を系統的に計算する:

\begin{center}
\small
\begin{tabular}{llllll}
\toprule
$q$ & $\chi$ & $L(-1,\chi)$ & $\Omega_\Lambda$ & 符号 & 物理的? \\
\midrule
2 & $\chi_0$ (主指標) & $+1/12$ & $2\pi/9 = 0.698$ & $+$ & \textbf{YES} \\
3 & $\chi_0$ (主指標) & $+1/6$ & $4\pi/9 = 1.396$ & $+$ & No ($>1$) \\
3 & $\chi_1$ (非主指標) & $0$ & $0$ & --- & 自明 \\
4 & $\chi_0$ (主指標) & $+1/12$ & $2\pi/9 = 0.698$ & $+$ & \textbf{YES}$^*$ \\
4 & $\chi_{-4}$ (奇) & $+1/4$ & $2\pi/3 = 2.094$ & $+$ & No ($>1$) \\
5 & $\chi_0$ (主指標) & $+1/3$ & $8\pi/9 = 2.793$ & $+$ & No ($>1$) \\
5 & その他 & 各種 & 各種 & & No \\
6 & $\chi_0$ (主指標) & $-1/6$ & $4\pi/9 = 1.396$ & $-$ & No (負) \\
6 & $\chi_{-3}$ & $+1/3$ & $8\pi/9 = 2.793$ & $+$ & No ($>1$) \\
\bottomrule
\end{tabular}
\end{center}

$^*$注: $\chi_0 \bmod 4$ は
$L(s,\chi_0) = \zeta(s)(1-2^{-s}) = \zeta_{\neg 2}(s)$ を与え，
$\chi_0 \bmod 2$ と同一である．$4 = 2^2$ の唯一の素因数が $2$ だからである．

\begin{theorem}[算術的ランドスケープの一意性]
\label{th:unique-landscape}
導手 $q \leq 100$ の全ての原始的Dirichlet指標のうち，
物理的真空 ($0 < \Omega_\Lambda < 1$) を与える $L$関数は
$k = 1, 2, 3, \ldots$ に対する主指標 $\chi_0 \bmod 2^k$ のみである．
これらは全て $L(s, \chi_0) = \zeta_{\neg 2}(s)$ を与え，
従って $\Omega_\Lambda = 2\pi/9$ となる．

算術的ランドスケープは\textbf{最大6個の真空} ($2^k \leq 64$) を含むが，
全て物理的に等価である．
\end{theorem}

\begin{proof}
主指標 $\chi_0 \bmod q$ に対して:
$L(-1, \chi_0) = \zeta(-1)\prod_{p|q}(1-p)$．

正値性は積 $\prod_{p|q}(1-p)$ が負であること
($\zeta(-1) < 0$ の符号を反転させるため) を要求する．
各因子 $(1-p) < 0$ であるから，積は $q$ の相異なる素因数の個数が
奇数のとき負になる．

大きさの制約:
$\Omega_\Lambda = (8\pi/3) \cdot (1/12) \cdot \prod_{p|q}|1-p| < 1$
は $\prod_{p|q}(p-1) < 9/(2\pi) \approx 1.432$ を要求する．

$q$ が $p \geq 3$ の奇素因数を持てば
$(p-1) \geq 2$ であるから $\prod(p-1) \geq 2 > 1.432$．不合格．

従って $q$ は $2$ のべき乗でなければならない: $q = 2^k$．
$q = 2^k$ に対して: 唯一の素因数は $2$ であり，$(2-1) = 1 < 1.432$．
因子数は $1$ (奇数)．符号は正．合格．

全ての $q = 2^k$ は $L(s, \chi_0) = \zeta(s)(1-2^{-s}) =
\zeta_{\neg 2}(s)$ を与え，$\Omega_\Lambda = 2\pi/9$ となる．

非主指標に対して: $L(-1, \chi)$ は零 (偶指標の $s = -1$ における自明な零点)
であるか，$q = 100$ までの全ての場合で $|\Omega_\Lambda(\chi)| > 1$ となる．
\end{proof}

\subsection{弦理論ランドスケープとの比較}

\begin{center}
\begin{tabular}{lll}
\toprule
& 弦理論ランドスケープ & 算術的ランドスケープ \\
\midrule
分類空間 & Calabi--Yau多様体 & Dirichlet $L$関数 \\
サイズ & $\sim 10^{500}$ & $\leq 6$ (実質的に1) \\
選択原理 & 人間原理 & 物理的制約 ($0 < \Omega < 1$) \\
$\Omega_\Lambda$ の決定 & 予測なし & $2\pi/9 = 0.698$ \\
反証可能性 & なし & あり (Euclid $0.5\%$) \\
\bottomrule
\end{tabular}
\end{center}

算術的ランドスケープは弦理論ランドスケープよりも単に``小さい''のではなく
--- \textbf{本質的に一意}である．
物理的制約は真空選択の曖昧さを許さない．

\section{情報論的解釈}

\subsection{境界エントロピーとAffleck--Ludwig $g$因子}

境界共形場理論 (BCFT) において，境界エントロピー
$s_{\rm bdy} = \log g$ は共形不変な境界条件を
特徴付ける~\cite{AffleckLudwig1991}．
自由ボソン ($c = 1$) に対して:

\begin{itemize}
\item Dirichlet境界: $g_D = 1/\sqrt{2}$
\item Neumann境界: $g_N = 1/\sqrt{2}$
\item 混合DN: $g_{DN} = 1$
\end{itemize}

$[0, L]$ 上の1次元場を $L/2$ で二分割した際の
エンタングルメントエントロピーは~\cite{CalabreseCardy2004}:
\begin{equation}
S = \frac{c}{6}\log\frac{L}{\epsilon} + \log g_{\rm bdy} + \ldots
\end{equation}

\subsection{DN境界の情報コスト}

DNとDDのエンタングルメントエントロピーの差:
\begin{align}
\Delta S &= S_{DN} - S_{DD} \\
&= \log g_{DN} - \log g_{DD} \\
&= \log 1 - \log(1/\sqrt{2}) \\
&= \frac{1}{2}\log 2
\end{align}

\begin{theorem}[ダークエネルギーの情報コスト]
\label{th:info-cost}
``標準的''DD真空から物理的DN真空への遷移は，
境界エンタングルメントエントロピーにおいて正確に
\begin{equation}
\boxed{\Delta S = \frac{1}{2}\log 2 \approx 0.347 \text{ nats}}
\end{equation}
のコストを要する．これは\textbf{半ビット}の情報量である．
\end{theorem}

\subsection{物理的解釈}

量 $\frac{1}{2}\log 2$ は物理学において以前から登場している:
\begin{itemize}
\item \textbf{Majoranaフェルミオン}: 境界上のMajoranaフェルミオンは
エンタングルメントエントロピーに $\frac{1}{2}\log 2$ を
寄与する~\cite{Kitaev2006}
\item \textbf{トポロジカルエンタングルメント}: 特定のトポロジカル相において，
普遍的な副次的項は $\gamma = \frac{1}{2}\log 2$ である
\item \textbf{$\mathbb{Z}/2$ ゲージ理論}: $\mathbb{Z}/2$ ゲージ理論の
トポロジカルエンタングルメントエントロピーは $\frac{1}{2}\log 2$ である
\end{itemize}

三つ全てが $\mathbb{Z}/2$ 構造を含む --- Euler積において $p = 2$ を選択し，
$K_1(\mathbb{Z}) = \mathbb{Z}/2$ を生成するのと同じ $\mathbb{Z}/2$ である．

\begin{observation}[$\mathbb{Z}/2$ 情報としてのダークエネルギー]
宇宙論定数は，宇宙のDN境界が $\frac{1}{2}\log 2$ nats の
境界エンタングルメントを担うことから生じる
--- $\mathbb{Z}/2$ トポロジカルセクターの情報論的シグネチャである．
ダークエネルギーは\emph{宇宙の時間的 $\mathbb{Z}/2$ 非対称性の
熱力学的コスト}である．
\end{observation}

\subsection{$\Delta S$ と $\Omega_\Lambda$ の接続}

境界エントロピー差 $\Delta S = \frac{1}{2}\log 2$ と
真空エネルギー係数 $|\zeta_{\neg 2}(-1)| = 1/12$ は
中心電荷 $c = 1$ を通じて関係する:

DN真空のCasimirエネルギー:
\begin{equation}
E_{DN} = -\frac{\pi c}{24 L}\cdot\zeta_{\neg 2}(-1)/\zeta(-1) = +\frac{\pi}{48 L}
\end{equation}

境界エントロピー:
\begin{equation}
\Delta S = \frac{1}{2}\log 2 = \log\frac{g_{DN}}{g_{DD}}
\end{equation}

両方の量はDN--DD差を符号化するが，異なるレベルにおいてである:
$\Delta S$ は\emph{情報}レベルで，
$\zeta_{\neg 2}(-1)$ は\emph{エネルギー}レベルで．

これらを結ぶ鎖:
\begin{equation}
\underbrace{\mathbb{Z}/2}_{\text{スピン/境界}}
\to \underbrace{p = 2}_{\text{Euler因子}}
\to \underbrace{\zeta_{\neg 2}(-1) = +\frac{1}{12}}_{\text{真空エネルギー}}
\to \underbrace{\Omega_\Lambda = \frac{2\pi}{9}}_{\text{ダークエネルギー}}
\end{equation}
\begin{equation}
\underbrace{\mathbb{Z}/2}_{\text{同一}}
\to \underbrace{\Delta S = \frac{1}{2}\log 2}_{\text{情報コスト}}
\to \underbrace{\text{境界エントロピー}}_{\text{Affleck--Ludwig}}
\end{equation}

両方の分岐は同じ $\mathbb{Z}/2$ から発する．

\section{算術的ランドスケープの一意性}

\subsection{三層選択}

物理的真空の一意性は，それぞれ異なる候補を排除する
三つの独立した制約から導かれる:

\textbf{第1層: 符号．} $\Omega_\Lambda > 0$ (宇宙の加速膨張) は
奇数個の素因数の除去を要求する．
これは偶数個の相異なる素因数を持つ全ての $q$
（例: $q = 6, 10, 14, 15, \ldots$）を排除する．

\textbf{第2層: 大きさ．} $\Omega_\Lambda < 1$ は
$\prod_{p|q}(p-1) < 9/(2\pi)$ を要求する．
$p \geq 3$ に対して $(p-1) \geq 2$ であるから，
奇素因数を持つ $q$ は全て不合格．$q = 2^k$ のみが生き残る．

\textbf{第3層: 等価性．} 生き残った $q = 2^k$ は全て
同一の $L$関数 $\zeta_{\neg 2}$ を与える．
ランドスケープは一点に崩壊する．

\begin{corollary}[真空縮退なし]
算術的ランドスケープには縮退がない: 全ての物理的真空は等価である．
この枠組みには``マルチバース''は存在しない．
\end{corollary}

\subsection{非主指標}

非主指標 $\chi \neq \chi_0$ に対して，$L$値 $L(-1, \chi)$ は
一般化Bernoulli数から計算される．
導手 $q = 100$ までの計算機による探索は，
全ての非主指標が $L(-1, \chi) = 0$
（偶指標の自明な零点）であるか
$|\Omega_\Lambda(\chi)| > 1$ を与えることを示す．

\begin{observation}
導手 $q \leq 100$ の非主指標Dirichlet指標で
物理的真空を与えるものは存在しない．
算術的ランドスケープは，素数べき導手 $2^k$ における
主指標によってのみ占められる．
\end{observation}

これは物理的解釈と整合する:
主指標 $\chi_0 \bmod q$ は $q$ を割る素数でのEuler因子を除去し，
これは真空モードの特定の境界条件に対応する．
非主指標は位相回転 $\chi(n) \neq 0, 1$ を含み，
明らかな境界条件の解釈を持たない．

\section{$\log 2$ の糸}

数 $\log 2 = 0.693\ldots$ はこの枠組みの至る所に
相互に関連した形で現れる:

\begin{enumerate}
\item \textbf{境界エントロピー}: $\Delta S = \frac{1}{2}\log 2$
（定理~\ref{th:info-cost}）

\item \textbf{$K_1$ レギュレータ}: 単元
$2 \in \mathbb{Z}[1/2]^\times$ のBorelレギュレータは
$\log|2| = \log 2$ である．これは
$K_1(\mathbb{Z}[1/2]) = \mathbb{Z}/2 \oplus \mathbb{Z}$
の基本レギュレータである~\cite{Milnor1971}

\item \textbf{Euler因子}: $(1 - 2^{-s})|_{s=0} = 1/2$ であり，
$\log(1/(1-2^{-s}))|_{s=0} = \log 2$

\item \textbf{自由エネルギー}: 境界自由エネルギー差
$\Delta F = T \Delta S = \frac{T}{2}\log 2$，
ここで $T$ はド・ジッター温度 $H/(2\pi)$

\item \textbf{Bekenstein限界}: 1ビットの情報は
エネルギー $k_BT\log 2$ に対応する．
半ビット: $\frac{1}{2}k_BT\log 2$
\end{enumerate}

$\log 2$ の5つの出現は全て同一の起源に遡る:
素数 $p = 2$ とそれに伴う $\mathbb{Z}/2$ 構造である．

\section{予測}

\begin{enumerate}
\item \textbf{$\Omega_\Lambda = 2\pi/9 = 0.698$}:
Euclid/DESIにより $0.5\%$ の精度で検証可能 (2025--2028)．
これは算術的ランドスケープの唯一の予測である．

\item \textbf{真空縮退なし}: 弦理論ランドスケープとは異なり，
異なる $\Lambda$ を持つ``ポケット宇宙''は存在しない．
マルチバースのシグネチャが発見された場合
（例: CMBにおけるバブル衝突），この枠組みは反証される．

\item \textbf{アナログ系における $\Delta S = \frac{1}{2}\log 2$}:
DN境界条件を持つ実験室系（例: トポロジカル絶縁体, Majoranaワイヤ）は
境界エントロピー $\frac{1}{2}\log 2$ を示すはずである．
これは既に凝縮系物理で観測されており，
独立した整合性の検証として機能する．

\item \textbf{ランドスケープの有限性}: 全てのDirichlet指標に対する
$L(-1, \chi)$ に関する将来の数学的研究により，
$\zeta_{\neg 2}$ が唯一の物理的真空であることが確認されるべきである．
これは数論的予測である．
\end{enumerate}

\section{限界}

\begin{enumerate}
\item 真空状態と $L$関数の同一視は\emph{提案}であり，
ラグランジアンから導出されたものではない．
``境界条件 $\leftrightarrow$ $L$関数'' の対応は
動機付けされているが証明されていない．

\item 式 (\ref{eq:omega-chi}) におけるCKN正規化 $(8\pi/3)$ は
ホログラフィック限界から仮定されている．
この正規化の第一原理からの導出は欠けている．

\item Affleck--Ludwig境界エントロピー $\frac{1}{2}\log 2$ は
$c = 1$ 自由ボソンに対して計算されたものである．
実際の宇宙論的場の内容はより豊かであり，
境界エントロピーは補正を受ける可能性がある．

\item $q = 100$ までの``非主指標探索''は計算機によるものであり，
一般的な証明ではない．全ての導手を覆う定理があれば
一意性の主張は強化される．

\item ダークマターおよびバリオン物質の密度は
この枠組みでは予測されない（確認された否定的
結果~\cite{wb-euler-sectors}）．
\end{enumerate}

\section{算術ホログラフィー：$p = 2$ はホログラフィック方向}

\subsection{コホモロジー次元と素数除去}

数論的スキームのエタールコホモロジー次元は，
空間の``幾何学的次元''の算術版として機能する．
Artin--Verdierの定理により:
\begin{equation}
\mathrm{cd}_{\text{\'et}}(\mathrm{Spec}(\mathbb{Z})) = 3
\end{equation}

一方，$p = 2$ を反転させた局所化 $\mathbb{Z}[1/2]$ に対して:
\begin{equation}
\mathrm{cd}_{\text{\'et}}(\mathrm{Spec}(\mathbb{Z}[1/2])) = 2
\end{equation}

\begin{observation}[ホログラフィック分解]
エタールコホモロジー次元において
$3 = 2 + 1$ という分解が成立する．
$\mathrm{Spec}(\mathbb{Z})$ の3次元構造は，
2次元の``境界'' $\mathrm{Spec}(\mathbb{Z}[1/2])$ と
1次元の $p = 2$ ``ホログラフィック方向''に分解される．
\end{observation}

この分解を時空の次元構造に拡張すると:
\begin{equation}
\underbrace{4}_{\text{時空}} =
\underbrace{2}_{\text{境界 } \mathrm{Spec}(\mathbb{Z}[1/2])}
+ \underbrace{1}_{p=2 \text{ ホログラフィック方向}}
+ \underbrace{1}_{\text{時間（WDW DN境界条件）}}
\end{equation}

DN真空（$p = 2$ の除去）は，この分解における
$p = 2$ ホログラフィック方向の``切除''に対応する．

\subsection{AdS/CFTとの接続}

反ド・ジッター/共形場理論（AdS/CFT）対応~\cite{Maldacena1999}は，
$d+1$ 次元の重力理論が $d$ 次元の場の理論と等価であることを主張する．
算術的ランドスケープにおける $p = 2$ のホログラフィック構造は，
AdS/CFTの算術版として解釈できる:

\begin{center}
\begin{tabular}{lll}
\toprule
& \textbf{AdS/CFT} & \textbf{算術版} \\
\midrule
バルク & AdS$_{d+1}$（$d+1$次元） & $\mathrm{Spec}(\mathbb{Z})$（$\mathrm{cd} = 3$）\\
境界 & CFT$_d$（$d$次元） & $\mathrm{Spec}(\mathbb{Z}[1/2])$（$\mathrm{cd} = 2$）\\
ホログラフィック方向 & 動径方向 $r$ & 素数 $p = 2$ の方向 \\
エンタングルメント & Ryu--Takayanagi公式 & $\Delta S = \frac{1}{2}\log 2$ \\
\bottomrule
\end{tabular}
\end{center}

特に，DN真空の境界エントロピー差
$\Delta S = \frac{1}{2}\log 2$ は，
\textbf{算術的Ryu--Takayanagi公式}として解釈できる:
$p = 2$ 方向に沿った``最小面''が境界エンタングルメントを決定し，
その値が半ビットとして現れる．

\subsection{遠アーベル的視点}

Neukirch--Uchidaの定理~\cite{NeukirchUchida}によれば，
絶対ガロア群 $\mathrm{Gal}(\bar{\mathbb{Q}}/\mathbb{Q})$ は
スキーム $\mathrm{Spec}(\mathbb{Z})$ を完全に決定する．
すなわち，代数的基本群が算術的``バルク時空''を復元する
--- これは遠アーベル的ホログラフィーの一形態である．

算術トポロジー~\cite{Morishita2012}の枠組みにおいて，
素数は結び目（knot）に対応する．
この類似の下で:
\begin{itemize}
\item 素数 $p$ $\leftrightarrow$ 3次元多様体中の結び目
\item $\mathrm{Spec}(\mathbb{Z}[1/p])$ $\leftrightarrow$ 結び目補空間（knot complement）
\item $p = 2$ でのDN真空 $\leftrightarrow$ $p = 2$ 結び目でのDehn手術
\end{itemize}

DN真空への遷移 --- $\zeta(s) \to \zeta_{\neg 2}(s)$ --- は，
$p = 2$ 結び目における\textbf{Dehn手術}（特定の手術係数によるトポロジー変更）
として理解される．
ダークエネルギー $\Omega_\Lambda = 2\pi/9$ は，
この算術的手術の\emph{物理的帰結}である:
$p = 2$ 結び目のDehn手術が時空のトポロジーを変え，
その変化が宇宙論定数として観測される．

\section{議論}

算術的ランドスケープは真空選択問題の根本的な単純化を表す．
弦理論が $10^{500}$ 個の真空と人間原理以外の選択原理を持たない一方で，
算術的ランドスケープは $\leq 6$ 個の等価な真空と
決定論的な選択 ($\Omega_\Lambda = 2\pi/9$) を提供する．

情報論的な糸は物理的深みを加える: ダークエネルギーは単なる宇宙論的
パラメータではなく，\emph{宇宙の $\mathbb{Z}/2$ 境界構造の
熱力学的コスト}である．$\frac{1}{2}\log 2$ の境界エントロピーは
宇宙論を凝縮系物理 (Majoranaフェルミオン, トポロジカルエンタングルメント)
および代数的 $K$理論 ($\mathbb{Z}[1/2]$ の $K_1$ レギュレータ) に接続する．

数論 ($L$関数)，情報理論 (境界エントロピー)，
代数的トポロジー ($K_1$, $\mathbb{Z}/2$)，
宇宙論 ($\Omega_\Lambda$) が
一つの数 $\log 2$ に収束することは，
我々の見解では偶然とは考えにくい．
素数 $p = 2$ が量子真空の算術的構造において
特別な役割を果たすことを示唆しており
--- その役割がダークエネルギーとして物理的に発現する．

この描像が正しいかどうかは，2028年までに期待される
Euclidによる $\Omega_\Lambda$ の $0.5\%$ 精度の測定によって決定される．

\section{結論}

我々は\textbf{算術的ランドスケープ}: Dirichlet $L$関数によって
分類される可能な真空状態の空間を提案した．
主要な結果は:

\begin{enumerate}
\item 物理的制約 ($0 < \Omega_\Lambda < 1$) が
ランドスケープを\textbf{ただ一つの真空}に帰着させる:
$\zeta_{\neg 2}(s) = (1-2^{-s})\zeta(s)$，
$\Omega_\Lambda = 2\pi/9$ を与える（定理~\ref{th:unique-landscape}）．

\item DN境界の (DDに対する) 情報コストは正確に
$\Delta S = \frac{1}{2}\log 2$ nats --- 半ビット ---
である（定理~\ref{th:info-cost}）．

\item ダークエネルギーは宇宙の時間的 $\mathbb{Z}/2$ 非対称性の
熱力学的コストであり，宇宙論を境界CFT，
トポロジカルエンタングルメント，$K$理論に接続する．

\item 算術的ランドスケープは $\leq 6$ 個の等価な真空を含み，
弦理論ランドスケープの $\sim 10^{500}$ に対する．
人間原理的選択は不要である．

\item $\mathrm{Spec}(\mathbb{Z})$ のエタールコホモロジー次元
$\mathrm{cd} = 3$ は，$\mathrm{Spec}(\mathbb{Z}[1/2])$ の
$\mathrm{cd} = 2$ と $p = 2$ ホログラフィック方向に分解される．
4次元時空は $2(\text{境界}) + 1(p=2) + 1(\text{時間})$ として
ホログラフィック的に構成され，DN真空は $p = 2$ 結び目での
算術的Dehn手術に対応する．
\end{enumerate}

\begin{thebibliography}{99}

\bibitem{Susskind2003}
L.~Susskind, hep-th/0302219.

\bibitem{KKLT2003}
S.~Kachru, R.~Kallosh, A.~Linde, S.~Trivedi,
Phys.\ Rev.\ D \textbf{68}, 046005 (2003).

\bibitem{Weinberg1987}
S.~Weinberg, Phys.\ Rev.\ Lett.\ \textbf{59}, 2607 (1987).

\bibitem{wb-functional}
Wright Brothers Research Program,
``Perturbative and non-perturbative wings of the zeta building''
(2026).

\bibitem{wb-wdw}
Wright Brothers Research Program,
``Wheeler--DeWitt temporal DN vacuum and $\Omega_\Lambda = 2\pi/9$''
(2026).

\bibitem{AffleckLudwig1991}
I.~Affleck and A.~W.~W.~Ludwig,
Phys.\ Rev.\ Lett.\ \textbf{67}, 161 (1991).

\bibitem{CalabreseCardy2004}
P.~Calabrese and J.~Cardy,
J.\ Stat.\ Mech.\ \textbf{0406}, P06002 (2004).

\bibitem{Kitaev2006}
A.~Kitaev and J.~Preskill,
Phys.\ Rev.\ Lett.\ \textbf{96}, 110404 (2006).

\bibitem{Milnor1971}
J.~Milnor, \textit{Introduction to Algebraic K-Theory}
(Princeton Univ.\ Press, 1971).

\bibitem{wb-euler-sectors}
Wright Brothers Research Program,
``Euler prime sectors: honest negative result'' (2026).

\bibitem{Maldacena1999}
J.~Maldacena,
Int.\ J.\ Theor.\ Phys.\ \textbf{38}, 1113 (1999);
hep-th/9711200.

\bibitem{NeukirchUchida}
J.~Neukirch, J.\ Reine Angew.\ Math.\ \textbf{302}, 147 (1978);
K.~Uchida, Ann.\ Math.\ \textbf{106}, 589 (1977).

\bibitem{Morishita2012}
M.~Morishita,
\textit{Knots and Primes: An Introduction to Arithmetic Topology}
(Springer, 2012).

\end{thebibliography}

\end{document}
